%%% FORMATO E IDIOMA %%%
    \documentclass[letterpaper,DIV=15,12pt]{scrartcl}
    \usepackage[spanish,mexico,shorthands=off,es-lcroman]{babel}
    \usepackage{scrlayer-scrpage}
        \clearpairofpagestyles
        \ihead{\footnotesize \textit{Teoría de Conjuntos III}}
        \ohead{\footnotesize \textit{2026-2}}
        \cfoot{\normalfont\thepage}
        \addtokomafont{title}{\bfseries \rmfamily}
        \setlength{\headsep}{5pt}
        \newpairofpagestyles{beginstyle}{
            \clearpairofpagestyles
            \KOMAoptions{headsepline=false}
            \cfoot{\footnotesize \pagemark}
            \cfoot{\normalfont\thepage}
        }
        \addtokomafont{section}{\raggedleft\rmfamily}
        \addtokomafont{subsection}{\raggedleft\rmfamily}
        \addtokomafont{subsubsection}{\raggedleft\rmfamily}
        \setlength{\headsep}{8pt}
        \setlength{\footskip}{20pt}
        \setlength{\textheight}{600pt}
%%% PAQUETES %%%
    \usepackage{ifthen}
    \usepackage[shortlabels]{enumitem}
        \setenumerate[1]{label=\MakeLowercase{\roman*}), ref=\roman*}
        \setenumerate[2]{label=\MakeLowercase{\alph*}), ref=\alph*}
    \usepackage{amsmath}
    \usepackage{amsthm}\usepackage[T1]{fontenc}
    \usepackage{stix2}
    \usepackage{stmaryrd}
	\renewcommand{\qedsymbol}{$\blacksquare$}
	\addto\captionsspanish{\renewcommand{\proofname}{\normalfont\bfseries Demostración}}	
	\newenvironment{solucion}{\renewcommand{\qedsymbol}{$\boxtimes$}\begin{proof}[\normalfont\bfseries Solución]}{\end{proof}}

    \usepackage{amsthm}
		%\newtheoremstyle{manual}{3pt}{3pt}{\normalfont}{}{\bfseries}{.}{.5em}{#1#3}
		%\theoremstyle{manual}
		\newtheorem{proposicion}{Proposición}
        \newtheorem{definicion}{Definición}
        \newtheorem{teorema}{Teorema}
        \newtheorem{lema}{Lema}
        \newtheorem{corolario}{Corolario}
        \newtheorem*{ejercicio*}{Ejercicio}        
%%% C O M A N D O S %%%
    \renewcommand{\emptyset}{\varnothing}
    \newcommand{\tq}{:}
    \newcommand{\yy}{\quad \text{ y } \quad}
    \newcommand{\sii}{\quad \text{ si y sólo si } \quad}
    \newcommand{\Bv}[1]{\llbracket #1 \rrbracket}
    \newcommand{\zfc}{\operatorname*{\textsc{Zfc}}}
%% D O C U M E N T O %%
\begin{document}
\thispagestyle{plain}

    \begin{center}
        {\fontsize{30}{60}\rmfamily \textbf{Teoremas de Frocing}} \\ \vspace{.2cm} Teoría de Conjuntos III, 2026-2
    \end{center}
    \begin{flushright}
        \footnotesize \hfill Profesor: Luis Jesús Turcio Cuevas.\\
        \hfill Ayudante: Hugo Víctor García Martínez.
    \end{flushright}
    
    Durante todo el documento, $M$ será un modelo estándar, transitivo (no necesariamente un ``modelo'' ETN) para $\zfc$. Este puede ser pensado como $V$, de hecho, nos olvidaremos de la nomenclatura ``$M$'' en cuanto apliquemos el método.

    \section*{$B$-nombres}

    Para cualquier álgebra de Boole $B$ se puede construir el universo booleano $M^B$. El lenguaje de forcing es un lenguaje que contiene al de la teoría de conjuntos, donde cada elemento de $M^B$ se agrega como constante (formal) nueva. Este es el motivo por el cual los elementos de $M^B$ se suelen denotar como $\dot{a}$ (esto es pensando en que, posteriormente, $a$ será la interpretación de la constante $\dot{a}$ en la extensión genérica).

    \begin{definicion}
        Los \textbf{$B$-nombres} serán los elementos del universo booleano $M^B$.

        El \textbf{lenguaje del forcing} es el lenguaje de primer orden con igualdad y tipo $\{ \in \} \cup M^{B}$, es decir, pensamos cada elemento de $M^B$ como una constante nueva de lenguaje.
    \end{definicion}

    Hay un tipo especial de nombres, los \textbf{nombres canónicos}, definidos para cada $x \in M$ de forma recursiva como a continuación se indica:
    \begin{equation}
        \check{x} = \{ (\check{y},1) \tq y \in x \} \, .
    \end{equation}

    Hay que tener un poco de cuidado, cada nombre canónico es claramente un nombre, pero el recíproco de esto es falso. Al momento de hacer el universo booleano $M^B$, pueden (y van a) aparecer nuevos nombres. La idea es que $M$ tiene una ``copia isomorfa'' dentro de $M^B$. Se tiene el siguiente resultado:
    \begin{proposicion}
        Sea $\varphi=\varphi(v_1, \dots, v_n)$ una fórmula del lenguaje de la teoría de conjuntos y supongamos que es $\Delta_0$. Entonces para cualesquiera conjuntos $a_1, \dots, a_n \in M$:
        \[ M \vDash \varphi(a_1, \dots, a_n) \sii \Bv{\varphi(\check{a_1}, \dots, \check{a_n})} = 1 \]
    \end{proposicion}

    Lo anterior tiene por consecuencia algo interesante ¿cuáles valores booleanos podría tomar una fórmula $\Delta_0$, $\varphi$, donde sus variables toman nombres canónicos?. Utilizando lo anterior, no es dificil deducir que sus valores booleanos sólo pueden ser $0$ o $1$, consecuentemente, para cualesquiera conjuntos $x,y \in M$,
    \[ x \in y \sii \Bv{\check{x} \in \check{y}} = 1 \, \text{ y} \]
    \[ x = y \sii \Bv{\check{x} = \check{y}} = 1 \, . \]
    En clase se probó algo más fuerte, pero de momento, nos servimos de esto para visualizar por qué hay una inclusión muy natural $M \hookrightarrow M^B$ (dada por $x \mapsto \check{x}$).
    
    \iffalse
    \newpage
    Con los nombres, pensados como constantes de lenguaje, consideramos su interpretación según el filtro $H$. Esta viene dada de manera recursiva, para cada $\dot{a} \in M^{\mathbb{P}}$, por:
    \begin{equation}
        \dot{a}^H = \{ \dot{t}^H \tq \exists p \in H \, (p \Vdash \dot{t} \in \dot{a}) \}
    \end{equation}

    \begin{definicion}
        La extensión genérica de $M$ bajo el filtro $H$ es el conjunto de todas las $H$-interpretaciones,
        \begin{equation}
            M[H] = \{ \dot{a}^H \tq \dot{a} \in M^{\mathbb{P}} \}  \, .
        \end{equation}
    \end{definicion}
    \fi

    \section*{Extensión genérica desde un álgebra de Boole}

    Cuando $U \subseteq B$ es cualquier ultrafiltro, podemos considerar los $B$-nombres, tratados como constantes de lenguaje e interpretarlos respecto a $U$ de la siguiente forma:
    \begin{equation}
        \dot{a}^U  = \{\dot{t}^U \tq \dot{a}(\dot{t}) \in U\}
    \end{equation}
    Definiendo la {\bfseries extensión genérica} como:
    \begin{equation}
        M[U]  = \{\dot{a}^U \tq \dot{a} \in M^B\}
    \end{equation}

    Con la extensión genérica creada de esta forma, se tienen los siguientes resultados.
    \begin{teorema}
        La extensión $M[U]$ es transitiva y para cada fórmula $\psi=\psi(v_1, \dots, v_n)$ ocurre:
        \[ M[U] \vDash \psi(\dot{a}^U_1 , \dots, \dot{a}^U _n) \sii \Bv{\psi(\dot{a_1}, \dots, \dot{a_n})} \in U \, . \]
    \end{teorema}
    
    Como en clase se demostró que cada nombre canónico $\check{x}$ se $U$-interpreta como $x$; y, que el valor booleano de cada axioma de $\operatorname*{\textsc{Zfc}}$ es $1$, se tiene lo posterior.
    \begin{enumerate}
        \item $M[U]$ es un modelo transitivo, estándar, para $\operatorname*{\textsc{Zfc}}$.
        \item $M \subseteq M[U]$.
    \end{enumerate}


    \section*{$\mathbb{P}$-nombres y relación de forcing}

    A partir de ahora, $\mathbb{P}=(P,\leq,1) \in M$ será un orden parcial en el modelo base, fijaremos también su completación booleana $e:P \to B^+$. 

    La idea es construir la extensión genérica y la relación de forcing, partiendo únicamente desde el orden parcial. Para ello convendremos que:
    \begin{enumerate}
        \item Un $\mathbb{P}$-nombre es un $B$-nombre.
        \item $M^{\mathbb{P}}=M^B$.
    \end{enumerate}

    El objetivo ahora es construir la relación de forcing, su utilidad primordial será la de poder hablar sobre lo que ocurre en la extensión genérica, pero desde dentro del modelo base.

    \begin{definicion}
        Dadas una fórmula en la teoría de conjuntos $\psi=\psi(v_1, \dots, v_n)$ y una condición $p \in P$; se define para cualesquiera nombres $\dot{a_1}, \dots, \dot{a_n} \in M^B$:
        \[ p \Vdash \varphi(\dot{a_1}, \dots, \dot{a_n} ) \sii e(p) \leq \Bv{\varphi(\dot{a_1}, \dots, \dot{a_n})} \, . \]
    \end{definicion}

    Nótese que bajo esta definición, tanto la completación booleana $(e,B)$ como la relación de forcing $\Vdash$ son definibles en el modelo base. Esto implica el \textbf{lema de definibilidad}; esto es, cada vez que $\psi=\psi(v_1, \dots, v_n)$ sea una fórmula de la teoría de conjuntos y $\dot{a_1}, \dots, \dot{a_n}$ sean $\mathbb{P}$-nombres, las colecciones del estilo:
    \[ \{ p \in P \tq p \Vdash \psi (\dot{a_1}, \dots, \dot{a_n}) \} \]
    serán conjuntos en el modelo base. Siempre que estas colecciones no involucren fórmulas con parámetros que sean un filtro genérico (o cualquier cosa que realmente \textit{se salga} del modelo base), serán conjuntos desde el punto de vista de $M$.

    \begin{teorema}
        La relación $\Vdash$ se comporta de la siguiente forma:
        \begin{enumerate}[a)]
            \item Monotonía.
                \begin{enumerate}[i)]
            \item Si $q \leq p$ y $p \Vdash \psi$, entonces $q \Vdash \psi$.
            \item Si $p \vDash \psi$ y $\Bv{\psi \to \alpha}=1$, entonces $p \Vdash \alpha$.
                \end{enumerate}
            \item Negación.
                \begin{enumerate}[i)]
            \item No existe $p \in P$ con $p \Vdash \psi$ y $p \Vdash \lnot \psi$ (simultáneamente).
            \item $p \Vdash \lnot \psi$ si y solamente si para cada $q \leq p$ se tiene $q \not\Vdash \psi$.
            \item $p \Vdash \psi$ si y sólo si $p \Vdash \lnot \lnot \psi$.
                \end{enumerate}
            \item Conjunciones.
                \begin{enumerate}[i)]
            \item $p \Vdash \psi \land \alpha$ si y sólo si $p \Vdash \psi$ y $p \Vdash \alpha$.
            \item $p \Vdash \forall x \, \psi(x)$ si y solamente si para cada $\dot{a} \in M^{\mathbb{P}}$ ocurre $p \Vdash \psi(\dot{a})$.
                \end{enumerate}
            \item Disyunciones.
                \begin{enumerate}[i)]
                    \item $p \Vdash \psi \lor \alpha$ si y sólo si $D=\{r \in P \tq r \Vdash \psi \text{ o } r \Vdash \alpha \}$ es denso bajo $p$.
                    \item $p \Vdash \exists x \, \psi (x)$ si y sólo si $D=\{r \in P \tq \exists \dot{a} \in M^{\mathbb{P}} \, ( r \Vdash \psi(\dot{a}) )\}$ es denso bajo $p$.
                \end{enumerate}
            \item Completez.
                \begin{enumerate}[i)]
                    \item Si $p \Vdash \exists x \, \psi(x)$, entonces existe $\dot{a} \in M^{\mathbb{P}}$ tal que $p \Vdash \psi(\dot{a})$.
                \end{enumerate}
      \end{enumerate}
    \end{teorema}
    \begin{proof}
        \textbf{(a)} Son inmediatas dada la definición de $\Vdash$.

        \textbf{(b)} Los incisos (i) y (iii) son inmediatos. Para (ii) supóngase primero que $p \Vdash \lnot \psi$. Si $q \leq p$, entones por (a.i) se sigue que $q \Vdash \lnot \psi$, obteniéndose de (b.i) que $q \not\Vdash \psi$.

        Ahora súpóngase que $p \not\Vdash \psi$, entonces $e(p) \not\leq \Bv{\psi}$ y consecuentemente $e(p) \Rightarrow \Bv{\psi} \neq 1$, de donde $e(p) \cdot \Bv{\psi}^* = e(p) \cdot \Bv{\psi} \neq 0$. Como $e$ tiene imagen densa, existe $s \in P$ de modo que $e(s) \leq e(p)$ y $e(s) \leq \Bv{\psi}$. Por lo anterior $e(s) = e(s) \cdot e(p) \neq 0$; como $e$ refleja compatibilidades, $s \parallel p$ en $\mathbb{P}$ y existe $q \in P$ con $q \leq p,s$. Así, $q \leq p$ y $e(q)\leq e(s) \leq \Bv{\psi}$, con lo cual, $q \Vdash \psi$.

        \textbf{(c)} Se realizará sólo el segundo inciso, el primero es muy similar, utiliza el hecho de que un elemento en $B$ es cota inferior de un conjunto si y sólo si es menor o igual que su ínfimo.

        Para (ii), basta notar que:
        \begin{align*}
            p \Vdash \forall x \, \psi(x) \sii & e(p) \leq \Bv{\forall x \, \psi(x)} \\
            \sii & e(p) \leq \prod_{\dot{a} \in M^{\mathbb{P}}} \Bv{\psi(\dot{a})} \\
            \sii & \text{Para cada } \dot{a}\in M^{\mathbb{P}}, \; e(p) \leq \Bv{\psi(\dot{a})} \\
            \sii & \text{Para cada } \dot{a}\in M^{\mathbb{P}}, \; p \Vdash \psi(\dot{a}) \, ,
        \end{align*}
        lo cual finaliza la prueba.

        \textbf{(d)} Utilizando (a.i), (b.ii) y (b.iii):
        \begin{align*}
            p \not\Vdash \alpha \lor \beta \sii & p \not\Vdash \lnot (\lnot \alpha \land \lnot \beta) \\
            \sii & \text{Existe } q \leq p \text{ tal que } q \Vdash \lnot \alpha \land \lnot \beta \\
            \sii & \text{Existe } q \leq p \text{ tal que } q \Vdash \lnot \alpha \text{ y } q \Vdash \lnot \beta \\
            \sii & \text{Existe } q \leq p \text{ tal que para cada } r \leq q \text{ ocurre } r \not\Vdash \alpha \text{ y } r \not\Vdash \beta \, ,
        \end{align*}
        probando que:
        \begin{align*}
            p \Vdash \alpha \lor \beta \sii & \text{Para cada } p \leq q \text{ existe } r \leq p \text{ tal que } r \Vdash \alpha \text{ o } r \Vdash \beta \\
            \sii & \text{Para cada } p \leq q \text{ existe } r \leq p \text{ tal que } r \in D \\
            \sii & D \text{ es denso bajo } p \, .
        \end{align*}

        El inciso (ii) se realiza de manera similar, utilizando que $\Bv{\lnot \exists x \, \psi(x) \Leftrightarrow \forall x \, \lnot \psi(x)}=1$.

        \textbf{(e)} Supóngase que $p \Vdash \exists x \, \psi(x)$. Recordando que $M^\mathbb{P}=M^B$, como modelo booleano-valuado, es completo, existe un nombre $\dot{a} \in M^B$ de manera tal que $\Bv{\exists x \, \psi(x)}=\Bv{\psi(\dot{a})}$. Resulta claro que entonces $p \Vdash \psi(\dot{a})$.
    \end{proof}

    \section*{Extensión genérica desde un orden parcial}

    A partir de ahora, $H \subseteq P$ será un filtro de $\mathbb{P}$, $M$-genérico y $G \subseteq B$ será el filtro generado por $H$ en $B$, es decir,
     \[ G=\uparrow e[H] = \{ b \in B \tq \exists p \in H \, ( e(p) \leq b) \} \, . \]

    Con los nombres, pensados como constantes de lenguaje, consideramos su interpretación según el filtro $H$. Esta viene dada de manera recursiva, similarmente a (2) y (3). Para cada $\dot{a} \in M^{\mathbb{P}}$ se define:
    \begin{equation}
        \dot{a}^H = \{ \dot{t}^H \tq \exists p \in H \, (e(p) \leq \dot{a}(\dot{t} ) \}
    \end{equation}

    Tomando ahora la extensión genérica de $M$, bajo el filtro $H$, como el conjunto de todos los $\mathbb{P}$-nombres ya interpretados,
        \begin{equation}
            M[H] = \{ \dot{a}^H \tq \dot{a} \in M^{\mathbb{P}} \}  \, .
        \end{equation}

    Para mostrar que en este método es lo mismo partir de un filtro genérico en un orden parcial, que en un filtro genérico de su completación booleana, bastará demostrar que $M[H]=M[G]$, para lo cual procederemos como a continuación se muestra.

    \begin{proposicion}
        Para cualquier nombre $a \in M^\mathbb{P} = M^B$, ocurre que $a^H = a^G$. Donde en la izquierda interpretamos como en (4), y a la derecha, como en (2). Consecuentemente,
        \[ M[G] = M[H] \, . \]
    \end{proposicion}
    \begin{proof}
        Por inducción supóngase que $\dot{a}$ es un $\mathbb{P}$- (o $B$-) nombre y que para cada elemento $x=\dot{t}^H \in \dot{a}^H$ ocurre $\dot{t}^H=\dot{t}^G$. Entonces, aplicando la hipótesis de inducción y la definición de $G$:
        \begin{align*}
            \dot{a}^G & = \{ \dot{t}^G \tq \dot{a}(\dot{t}) \in G \} \\
            & = \{ \dot{t}^G \tq \exists p \in H \, (e(p) \leq \dot{a}(\dot{t})) \\
            & = \dot{a}^H \, ,
        \end{align*}
        finalizando la prueba.
    \end{proof}

\end{document}
